\documentclass[12pt,a4paper]{article}
\usepackage[utf8]{inputenc}
\usepackage[T1]{fontenc}
\usepackage[spanish]{babel}
\usepackage[margin=1.4cm]{geometry}
\usepackage{array,tabularx}
\usepackage{enumitem}
\usepackage{titlesec}
\usepackage{xcolor}
\usepackage{hyperref}

% Restrained, high-contrast palette shared by headings, emphasis and links.
\definecolor{HeadingColor}{HTML}{243B53}
\definecolor{InstructionColor}{HTML}{704214}
\definecolor{ConceptColor}{HTML}{205C50}
\definecolor{LinkColor}{HTML}{0047AB}
\hypersetup{colorlinks=true,allcolors=LinkColor,pdfborder={0 0 0}}
\newcommand{\guidelink}[2]{\href{#1}{\textcolor{LinkColor}{\textbf{\underline{#2}}}}}
\newcommand{\instruction}[1]{\textcolor{InstructionColor}{\textbf{#1}}}
\newcommand{\concept}[1]{\textcolor{ConceptColor}{\textbf{#1}}}

\setlength{\parindent}{0pt}
\setlength{\parskip}{3pt}
\setlength{\tabcolsep}{5pt}
\renewcommand{\arraystretch}{1.12}
\setlist{leftmargin=0.6cm,topsep=3pt,partopsep=0pt,parsep=0pt,itemsep=1pt}
\titleformat{\section}{\normalsize\bfseries\color{HeadingColor}}{}{0pt}{}
\titlespacing*{\section}{0pt}{8pt}{3pt}
\newcolumntype{Y}{>{\raggedright\arraybackslash}X}
\newcommand{\checkbox}{\fbox{\rule{0pt}{0.8ex}\rule{0.8ex}{0pt}}}

\begin{document}
\begin{center}
{\Large\bfseries\color{HeadingColor} C9--C10: Explorar ideas y comprender al usuario.}\\[3pt]
Emprendimiento e Innovación --- Grado 10
\end{center}

\section*{Objetivo y punto de partida}
\concept{C9:} Participar en actividades iniciales de ideación con apertura a propuestas diversas.\\
 \concept{C10:} Reconocer la importancia de comprender al usuario antes de proponer una solución. 
 
C5 organizó la evidencia; C6 priorizó una oportunidad; C7 definió el desafío y C8 conectó conocimientos. Ahora generarán posibilidades y examinarán cómo la comprensión del usuario orienta sus ideas. \instruction{Las ideas son exploratorias: esta etapa termina sin elegir, desarrollar, prototipar ni implementar una solución final.}

\section*{Instrucciones de trabajo y entrega}
\instruction{Esta guía es de consulta.} Continúen el mismo documento acumulativo del proyecto en LaTeX: conserven las secciones anteriores y añadan únicamente las nuevas. Mantengan el título y actualicen fecha e integrantes si cambian. Entreguen el PDF actualizado; pueden usar \guidelink{https://www.overleaf.com/}{Overleaf}. No repitan perfiles, mapas, matrices ni explicaciones: remitan a los códigos y secciones existentes.

La IA es opcional y solo apoya el código, el formato y la organización de contenido propio. \instruction{No debe inventar evidencia, conclusiones, decisiones, ideas ni aportes de integrantes.} Distingan observaciones, conclusiones, supuestos e ideas propuestas. Los ejemplos son ficticios, no resultados del grupo.
\section*{C9 --- 1. Recuperen el desafío y sus puntos de apoyo}
Lean D1, la representación del usuario y los hallazgos H de C5, y K1/K2 de C8. Antes de idear, comprueben qué necesidad está respaldada y qué sigue siendo un supuesto. En la nueva sección incluyan solo referencias a estos resultados; no copien la pregunta ni escriban otro resumen. Si detectan una inconsistencia, corrijan la sección de origen y anoten el cambio y su razón.
\section*{C9 --- 2. Generen ideas de manera individual}
Durante cinco minutos, cada integrante escriba \textbf{tres ideas diferentes} para responder al desafío, sin consultar ni evaluar las propuestas de otros. Describan cada idea en una frase: qué podría hacerse y cómo intentaría atender la necesidad. Busquen distintas maneras de actuar, organizar, comunicar o facilitar la experiencia; cambiar solo el nombre de una propuesta no crea una idea diferente.

Identifiquen autor y código de cada idea (I1, I2, etc.). Conserven todas para compartirlas. No necesitan dibujos, presupuestos ni detalles técnicos.

\section*{C9 --- 3. Escuchen, combinen y amplíen}
\begin{enumerate}
\item Por turnos, cada integrante presenta sus tres propuestas. Los demás escuchan sin interrumpir ni descartarlas de inmediato. Antes de opinar, hagan una pregunta para comprender una propuesta ajena.
\item Cada integrante aporte una manera de ampliar, adaptar o combinar una idea de otra persona. Pueden apoyarse en los conocimientos de C8 para mirarla desde otra perspectiva.
\item Registren \textbf{dos propuestas nuevas o modificadas del grupo}, indicando las ideas de origen y qué cambió. No voten ni ordenen las ideas de mejor a peor: todavía exploran posibilidades.
\end{enumerate}
Dejen un registro breve de participación: \textbf{integrante, idea ajena que escuchó, pregunta realizada y aporte}. Este registro debe mostrar cómo consideraron una perspectiva diferente, no solo quién estuvo presente. Si una propuesta les parece extraña, pregunten qué necesidad intenta atender antes de juzgarla.

\section*{C9 --- 4. Organicen el conjunto de ideas}
Reúnan \textbf{todas las ideas individuales y las dos propuestas colaborativas} en una sola tabla. Una fila por idea; no repitan las descripciones en otro listado.

{\small
\begin{tabularx}{\linewidth}{|p{0.08\linewidth}|Y|Y|Y|}
\hline
\textbf{Código} & \textbf{Idea: descripción breve} & \textbf{Necesidad que intenta atender} & \textbf{Autoría y origen} \\
\hline
I1, etc. & Qué proponen explorar. & Necesidad del usuario, con referencia a C5. & Integrante; o grupo, códigos de origen y modificación. \\
\hline
\end{tabularx}
}

Usen los códigos de origen para mostrar las combinaciones en la misma tabla. \concept{Una idea inicial es una posibilidad por explorar, no una respuesta demostrada.}

\section*{C10 --- 5. Contrasten dos ideas con lo que saben del usuario}
Revisen la representación del usuario y los hallazgos referenciados en la actividad 1. Tomen \textbf{dos ideas distintas} de la tabla, al menos una surgida del trabajo colaborativo, para examinar su relación con los usuarios. Esta muestra sirve para reflexionar, no para elegir finalistas. Para cada una, completen:

{\small
\begin{tabularx}{\linewidth}{|p{0.08\linewidth}|Y|Y|Y|}
\hline
\textbf{Idea} & \textbf{¿Qué sabemos del usuario?} & \textbf{¿Qué estamos suponiendo?} & \textbf{¿Qué falta investigar?} \\
\hline
Código. & Necesidad o dificultad observada; hallazgo y registro que la respaldan. & Qué tendría que ser cierto para que la idea respondiera a esa necesidad. & Pregunta pendiente y cómo podrían aclararla más adelante. \\
\hline
\end{tabularx}
}

La evidencia de una necesidad \textbf{no demuestra que una idea funcionará ni que el usuario la aceptará}. Si no encuentran respaldo, indíquenlo: no inventen una observación para justificar la propuesta. Añadan a cada fila una condición que la idea debe respetar según la evidencia (por ejemplo, no exigir recursos a los que el usuario no tiene acceso). Si requiere ajuste, actualicen la idea en su fila original y anoten aquí su código, el cambio y el hallazgo que lo motiva. Si la mantienen, justifiquen qué condición respeta. No asignen puntajes ni comparen ganadoras. \instruction{Aquí identifican vacíos; no deben realizar nuevas entrevistas ni observaciones.}

\section*{C10 --- 6. Expliquen por qué comprender al usuario orienta las ideas}
Escriban una conclusión de tres o cuatro frases referida a una idea examinada: ¿qué riesgo tendría proponerla basándose solo en gustos personales?, ¿cómo la evidencia del usuario obliga a ajustarla o limitarla?, ¿qué supuesto deben verificar antes de desarrollarla? No repitan la tabla. La necesidad documentada orienta la propuesta, pero no demuestra su aceptación ni su funcionamiento.

\section*{Resultado y continuidad}
C9 aporta un registro único de ideas con autoría y evidencia de escucha y colaboración. C10 aporta el contraste de dos ideas, sus condiciones y una explicación fundamentada de por qué comprender al usuario debe orientar las propuestas. Conserven las posibilidades abiertas. Los pendientes informarán una futura investigación; esta guía no pide seleccionar soluciones, prototipar ni implementar.
\section*{Verificación final}
{\small
\begin{itemize}[label=\checkbox]
\item Recuperamos el desafío de C7, los hallazgos de C5 y los conocimientos de C8 sin repetir sus actividades.
\item Cada integrante generó tres ideas propias y dejó evidencia de escucha y de un aporte a una propuesta ajena.
\item Documentamos todas las ideas y dos combinaciones o modificaciones, con autoría, origen y necesidad atendida.
\item Contrastamos dos ideas con registros reales y distinguimos lo que sabemos, suponemos y falta investigar.
\item El registro muestra apertura a otras propuestas; C10 explica cómo la evidencia del usuario condiciona las ideas y qué falta verificar.
\item La IA no inventó aportes ni evidencia; compilamos y revisamos el PDF acumulativo.
\end{itemize}
}
\end{document}